\section{Результаты}
\label{sec:results}

Все величины ниже получены в одном сквозном расчёте при фиксированных входных
данных раздела~\ref{sec:methods}; каждая сопровождается собственной проверкой
сходимости с объявленным допуском.

\subsection{Дисперсионная модель золота}
\label{sec:res-material}

Качество рациональной модели~\eqref{eq:lorentz-fit} измеряется относительно
табличной поляризуемости, вычисленной по~\eqref{eq:alpha-L} из данных
Джонсона--Кристи. Использованы две нормированные меры,
\begin{equation}
  \mathrm{NRMS}(X) =
  \frac{\bigl\langle |X_{\rm fit} - X_{\rm tab}|^2 \bigr\rangle^{1/2}}
       {\bigl\langle |X_{\rm tab}|^2 \bigr\rangle^{1/2}},
  \qquad
  \epsilon_{\max}(X) =
  \frac{\max_E |X_{\rm fit} - X_{\rm tab}|}{\max_E |X_{\rm tab}|},
  \label{eq:fit-metrics}
\end{equation}
применённые и к $X=\alpha$, и к $X=1/\alpha$: обратная величина входит в
знаменатели связанных уравнений, поэтому её точность контролируется отдельно.

\begin{table}[t]
  \centering
  \caption{Точность дисперсионной модели в окне подгонки
  $0.8$--$\SI{3.5}{\electronvolt}$ (рис.~1).}
  \label{tab:material}
  \begin{tabular}{llcccc}
    \toprule
    & & \multicolumn{2}{c}{$N=1$} & \multicolumn{2}{c}{$N=12$} \\
    \cmidrule(lr){3-4}\cmidrule(lr){5-6}
    Ориентация & Величина & NRMS & $\epsilon_{\max}$ & NRMS & $\epsilon_{\max}$ \\
    \midrule
    продольная & $\alpha$   & 0.722 & 0.774 & 0.0248 & 0.0246 \\
    продольная & $1/\alpha$ & 0.528 & 0.549 & 0.0248 & 0.0479 \\
    поперечная & $\alpha$   & 0.431 & 0.536 & 0.0238 & 0.0408 \\
    поперечная & $1/\alpha$ & 0.365 & 0.450 & 0.0181 & 0.0456 \\
    \bottomrule
  \end{tabular}
\end{table}

Одноосцилляторная модель ошибается на десятки процентов
(табл.~\ref{tab:material}): её NRMS составляет $0.722$ для продольной
поляризуемости, тогда как двенадцатиполюсная даёт $0.0248$ — выигрыш в $29$
раз. Это ожидаемо: в окне $0.8$--$\SI{3.5}{\electronvolt}$ золото имеет и
друдевскую часть, и межзонные переходы, которые один лоренцев контур
воспроизвести не может. Ниже (раздел~\ref{sec:res-material-transfer}) показано,
во что эта разница обходится в нелинейном режиме.

\subsection{Слабополевой отклик: усиление, сдвиг, уширение}
\label{sec:res-weak}

Рис.~2a--2c построены по трём величинам, определённым
в~\eqref{eq:gain}, и по сдвигу и ширине спектра~\eqref{eq:S}, нормированным на
ширину линии изолированной КТ $\Gamma_0=\SI{2.541}{\milli\electronvolt}$:
\begin{equation}
  \frac{\Delta E_{\rm res}(g)}{\Gamma_0}
  = \frac{E_{\rm res}^{\rm hyb}(g) - E_{\rm res}^{\rm iso}}{\Gamma_0},
  \qquad
  \frac{\Gamma_{\rm eff}(g)}{\Gamma_0}.
  \label{eq:shift-width}
\end{equation}

Знак эффекта задаётся не расстоянием, а ориентацией поля относительно
поверхности частицы в точке расположения КТ. Для двух каналов, где поле
нормально к поверхности, отклик усиливается; для трёх касательных —
подавляется. При $g=\SI{1}{\nano\meter}$ значения $G_{\rm exc}$ составляют
$25.3$ и $4.73$ против $0.066$, $2.19$ и $0.013$ соответственно.

Зависимость от зазора для ведущего канала (КТ у вершины, поле вдоль оси)
немонотонна:
\begin{center}\small
\begin{tabular}{lcccccc}
  \toprule
  $g$, нм & 0.5 & 0.75 & 1.0 & 1.5 & 2.0 & 3.0 \\
  \midrule
  $G_{\rm exc}$ & 20.2 & 23.5 & 25.3 & \textbf{25.5} & 23.4 & 18.2 \\
  $\tilde G_{\rm exc}$ & \textbf{38.0} & 35.9 & 33.6 & 29.0 & 24.9 & 18.5 \\
  $\Delta E_{\rm res}/\Gamma_0$ & $-0.544$ & $-0.411$ & $-0.318$ & $-0.201$ & $-0.135$ & $-0.068$ \\
  $\Gamma_{\rm eff}/\Gamma_0$ & 1.155 & 1.125 & 1.104 & 1.074 & 1.055 & 1.032 \\
  \bottomrule
\end{tabular}
\end{center}

Усиление на фиксированной частоте имеет \emph{внутренний} максимум вблизи
$1.25$--$\SI{1.5}{\nano\meter}$, тогда как усиление с подстройкой лазера растёт
монотонно вплоть до наименьшего просчитанного зазора. Причина видна в той же
таблице: сближение одновременно усиливает ближнее поле и смещает гибридный
резонанс, а при $g=\SI{0.5}{\nano\meter}$ смещение $-\SI{1.38}{\milli\electronvolt}$
составляет уже половину ширины линии $\SI{2.94}{\milli\electronvolt}$. Узкополосный
зонд на неподвижной частоте теряет на расстройке больше, чем выигрывает на поле;
перестраиваемый зонд не теряет ничего. Уширение остаётся умеренным: не более
$16\,\%$ во всём диапазоне.

\subsection{Расхождение дипольной и точной моделей по спектру}
\label{sec:res-discrepancy}

Мера расхождения на фиксированном окне признака $W$ (рис.~2d):
\begin{equation}
  \delta_{\rm spec}(g) =
  \left[\frac{\int_W \bigl|S_{\rm DD}(E) - S_{\rm FQS}(E)\bigr|^2\,dE}
             {\int_W S_{\rm FQS}^2(E)\,dE}\right]^{1/2}.
  \label{eq:delta-spec}
\end{equation}
Ни один из спектров не нормируется отдельно, поэтому~\eqref{eq:delta-spec}
чувствительна и к форме, и к амплитуде. Допуск $\delta_{\rm spec}\le0.1$
объявлен до расчёта.

Расхождение монотонно убывает с расстоянием, и в дальней зоне
($g\ge\SI{3}{\nano\meter}$) описывается степенным законом $\delta_{\rm
spec}\propto R^{-p}$. Для ведущего канала $p=2.97$ при максимальной
относительной невязке подгонки $3.5\,\%$, то есть с хорошей точностью
$\delta_{\rm spec}\propto R^{-3}$.

Тем не менее допуск в допустимой области не достигается. Для трёх слабых
каналов граница применимости DD установлена (зазоры $7$--$\SI{10}{\nano\meter}$),
для двух каналов с наибольшим усилением — нет:
$\delta_{\rm spec}$ ведущего канала падает лишь с $0.694$ при
$g=\SI{0.5}{\nano\meter}$ до $0.157$ при $g=\SI{12}{\nano\meter}$.

Экстраполяция измеренного закона объясняет, почему. Допуск $0.1$ был бы
достигнут при $R\approx\SI{34}{\nano\meter}$, где $k_m R\approx0.58$, тогда как
квазистатическое условие $k_m R\le0.5$ ограничивает расстояние величиной
$R\le\SI{29.4}{\nano\meter}$. Два требования разминулись примерно на $15\,\%$ по
расстоянию:
\begin{equation}
  R_{\delta\le0.1} \approx \SI{34}{\nano\meter}
  \;>\;
  R_{k_mR\le0.5} = \SI{29.4}{\nano\meter}.
  \label{eq:no-window}
\end{equation}
То есть дипольное приближение достигает объявленной точности только там, где уже
отказывает сама квазистатика, и ни один зазор не удовлетворяет обоим условиям
одновременно.

Экстраполяция проверяема: та же процедура, применённая к трём каналам, где
граница \emph{была} установлена, предсказывает зазоры $9.0$, $9.2$ и
$\SI{6.9}{\nano\meter}$ против найденных в расчёте $10$, $10$ и
$\SI{7}{\nano\meter}$ — в каждом случае это ближайший узел сетки не ниже
предсказания. Для пятого канала (КТ сбоку, поле вдоль оси) степенной закон
данные не описывает, и экстраполяция для него не приводится.

\subsection{Порог нелинейного возбуждения}
\label{sec:res-threshold}

Основная наблюдаемая — порог~\eqref{eq:threshold}, приведённый к порогу
изолированной КТ $\mathcal F_{\eta,0}=\SI{6.167e-5}{\joule\per\centi\meter\squared}$:
\begin{equation}
  \mathcal{R}(g) = \frac{\mathcal F_\eta(g)}{\mathcal F_{\eta,0}},
  \qquad
  \delta_{\mathcal F}(g) =
  \frac{\bigl|\mathcal F^{\rm DD}_\eta(g) - \mathcal F^{\rm FQS}_\eta(g)\bigr|}
       {\mathcal F^{\rm FQS}_\eta(g)}.
  \label{eq:threshold-metrics}
\end{equation}
Величина $\mathcal{R}<1$ означает выигрыш; $\delta_{\mathcal F}$ (рис.~3b)
переносит сравнение моделей с линейного спектра на центральную нелинейную
метрику работы.

\begin{table}[t]
  \centering
  \caption{Порог $\mathcal F_\eta$ в точной модели, $10^{-6}$~Дж/см$^2$
  (рис.~3a). Прочерк — порог не достигнут во всём сканируемом диапазоне, то
  есть частица \emph{ухудшает} возбуждение.}
  \label{tab:thresholds}
  \small
  \begin{tabular}{lcccccc}
    \toprule
    Канал & 0.5 нм & 1.0 нм & 2.0 нм & 5.0 нм & 12 нм \\
    \midrule
    вершина, вдоль оси          & \textbf{1.35} & 1.62 & 2.32 & 5.56 & 17.9 \\
    бок, поперёк радиально      & 8.50 & 9.39 & 11.5 & 19.0 & 35.7 \\
    бок, вдоль оси              & 21.1 & 26.9 & 45.3 & ---  & ---  \\
    вершина, поперёк            & ---  & ---  & ---  & 136  & 79.9 \\
    бок, поперёк касательно     & ---  & ---  & ---  & 200  & 88.8 \\
    \bottomrule
  \end{tabular}
\end{table}

Наилучшая конфигурация — КТ у вершины при поле вдоль оси
(табл.~\ref{tab:thresholds}): порог снижается в
\begin{equation}
  \mathcal{R}^{-1} = \frac{6.167\cdot10^{-5}}{1.350\cdot10^{-6}} = 45.7
  \label{eq:best-ratio}
\end{equation}
раз относительно изолированной КТ. Разброс по каналам при этом шире, чем по
зазору: выбор канала меняет результат от выигрыша в $45.7$ раза до полного
отсутствия порога, тогда как весь диапазон зазоров внутри лучшего канала даёт
множитель $13.3$.

Профиль порога по зазору насыщается. Между $3$ и $\SI{1}{\nano\meter}$ он
улучшается вдвое, а от $1$ нм до наименьшего просчитанного зазора — лишь в
\begin{equation}
  \frac{\mathcal F_\eta(\SI{1.0}{\nano\meter})}{\mathcal F_\eta(\SI{0.5}{\nano\meter})}
  = 1.20
  \label{eq:saturation}
\end{equation}
раза. Как следствие, пороги при $0.5$ и $\SI{0.75}{\nano\meter}$
($1.350$ и $1.478\cdot10^{-6}$~Дж/см$^2$, различие $9.5\,\%$) неразличимы в
пределах объявленной численной разрешающей способности $5\,\%$, применяемой к
обоим значениям. Поэтому результат сформулирован как интервал
$g\in[0.5,\,0.75]\ \text{нм}$, вне которого все конфигурации надёжно хуже, а не
как одна точка.

Расхождение моделей по порогу ведёт себя так же, как по спектру, но оно
систематически больше: для ведущего канала $\delta_{\mathcal F}$ убывает с
$1.53$ при $g=\SI{0.5}{\nano\meter}$ до $0.186$ при $\SI{12}{\nano\meter}$, ни
разу не опускаясь ниже допуска. При $g=\SI{1}{\nano\meter}$ дипольная модель
завышает требуемый флюенс в $2.4$ раза ($3.875$ против
$1.620\cdot10^{-6}$~Дж/см$^2$). Существенно, что она ошибается и качественно:
для канала «вершина, поперёк» DD даёт конечный порог при всех зазорах, тогда как
точная модель показывает, что при $g\le\SI{3}{\nano\meter}$ порог недостижим.

\subsection{Перенос материальной ошибки в наблюдаемые}
\label{sec:res-material-transfer}

Рис.~4 показывает спектр~\eqref{eq:S} в выбранной геометрии для трёх
материальных представлений, рис.~5 — нелинейную зависимость $P_{\rm exc}(\mathcal
F)$ и соответствующие пороги. Мерой переноса служит отношение порогов,
полученных при одинаковой геометрии и различающихся только числом полюсов:
\begin{equation}
  \mathcal{Q}(g) =
  \frac{\mathcal F_\eta\bigl(N=1\bigr)}{\mathcal F_\eta\bigl(N=12\bigr)}.
  \label{eq:material-ratio}
\end{equation}

Для ведущего канала $\mathcal{Q}=7.74$: одноосцилляторная модель завышает
требуемый флюенс почти на порядок ($1.045\cdot10^{-5}$ против
$1.350\cdot10^{-6}$~Дж/см$^2$). Для двух других разрешённых каналов
$\mathcal{Q}=4.59$ и $1.90$. Максимальное расхождение населённостей на общей
сетке флюенсов достигает $|{\Delta P}|=0.868$, то есть почти всего доступного
диапазона.

Таким образом, различие в табл.~\ref{tab:material}, выглядящее как «$72\,\%$
против $2.5\,\%$» в линейной характеристике материала, переносится в нелинейный
режим как ошибка порога в $7.7$ раза. Одного лоренцева осциллятора для золота в
этой задаче недостаточно.

\subsection{Временна\char'27{}я динамика}
\label{sec:res-dynamics}

Рис.~6 сравнивает $\rho_{ee}(t)$ для изолированной КТ, точной и дипольной
моделей при \emph{общем} импульсе и общем моменте считывания. В ближнем случае
($g=\SI{0.5}{\nano\meter}$, флюенс, отвечающий порогу точной модели) населённость
к моменту считывания равна $0.500$ в точной модели и $0.226$ в дипольной: при
одном и том же импульсе DD недооценивает возбуждение в $2.2$ раза. Контрольный
канал даёт $4.1\cdot10^{-4}$ — на три порядка ниже, что и есть количественная
мера подавления при касательном поле.

В дальнем случае ($g=\SI{12}{\nano\meter}$) те же величины составляют $0.0460$ и
$0.0388$: модели сблизились, но расхождение $19\,\%$ сохраняется, что
согласуется с $\delta_{\mathcal F}=0.186$ из
раздела~\ref{sec:res-threshold}.

\subsection{Устойчивость результата}
\label{sec:res-stability}

Ранжирование конфигураций проверено против одиночных вариаций входных
параметров. Относительное изменение порога составляет: зазор $\pm\SI{0.2}{\nano\meter}$
— $7$--$8\,\%$; полуоси $\pm5\,\%$ — $22$--$33\,\%$; энергия перехода
$\pm\SI{5}{\milli\electronvolt}$ — $3$--$6\,\%$; $\gamma_1\pm20\,\%$ —
$0.04\,\%$; $\Gamma_2\pm20\,\%$ — $0.6$--$1.9\,\%$. Во всех проверенных случаях
лучший канал и лучшая конфигурация не меняются.

Численные уточнения дают изменения много меньше объявленных допусков:
удвоение пространственного порядка — $3.6\cdot10^{-15}$, удвоение сетки
флюенса — $4.0\cdot10^{-6}$, уменьшение шага по времени — $3.6\cdot10^{-7}$,
переход к $N=13$ полюсам — $2.9\cdot10^{-4}$. Пороги, полученные двумя
независимыми путями (карта рис.~3 и материальная серия рис.~5), совпадают с
относительной разницей $2.8\cdot10^{-6}$.

Отдельно отмечено, что при несущей, смещённой на $+\SI{40}{\milli\electronvolt}$,
ведущая конфигурация не достигает $\eta=0.5$ на первой ветви Раби, и
ранжирование при такой несущей не определено. Это не расходимость расчёта, а
свойство задачи: отстройка лазера — управляемый параметр, и рабочая точка должна
выбираться вблизи резонанса.
